% =================================================================================================
% TEMPLATE for a results report on one or more production runs.  UNFINISHED BY DESIGN.
%
% Read Report/README.md first. It says when to write a report, where every number must come from,
% and the checklist to clear before the PDF goes to anyone. The reference example of a finished
% report is Report/populations/populations_report.tex -- when in doubt, do what it does.
%
% Each report lives in its own directory, Report/<topic>/, with its build products beside it.
% The template, the shared refs.bib and the README stay in Report/ itself. To start a report:
%     mkdir Report/<topic> && cp Report/TEMPLATE_report.tex Report/<topic>/<topic>_report.tex
%     edit the \graphicspath line, fill every \TODO, delete sections that do not apply
%     cd Report/<topic> && latexmk -pdf <topic>_report.tex
%     grep -c '^TODO:' <topic>_report.log     # must be 0 before the report is shared
% The paths below (../../figures, ../refs) assume exactly that depth.
%
% Every \TODO prints in red in the PDF AND writes a line to the .log, so an unfinished report
% cannot be mistaken for a finished one. The %% comments are instructions; delete them as you go.
% =================================================================================================
\documentclass[english,10pt,a4paper]{article}
\usepackage[T1]{fontenc}
\usepackage[margin=2.4cm]{geometry}
\usepackage{graphicx}
%% Point at the directories the analysis scripts WROTE the figures to. Never copy figures into
%% Report/: a copy drifts from the script that made it, the original cannot.
\graphicspath{{./}{../../figures/TODO_run_dir/}}
\usepackage{mathtools}
\usepackage{amssymb}
\usepackage{booktabs}
\usepackage{xcolor}
\usepackage[colorlinks=true,allcolors=blue]{hyperref}
\usepackage{caption}
\usepackage[authoryear,round]{natbib}
%% Keeps floats inside their own section; without it tables land pages after their text.
\usepackage[section]{placeins}

\newcommand{\tE}{t_{\rm E}}
\newcommand{\piE}{\pi_{\rm E}}
\newcommand{\tetE}{\theta_{\rm E}}
\newcommand{\Ml}{M_{\rm L}}
\newcommand{\Msun}{M_\odot}
\newcommand{\Neff}{N_{\rm eff}}

%% Placeholder: red in the PDF, and logged so `grep '^TODO:' *.log` finds every one left.
\newcommand{\TODO}[1]{\textcolor{red}{\textbf{[TODO: #1]}}\typeout{TODO: #1}}

\title{\TODO{title: what was run, and the one question it answers}}
\author{Ali Salehi}
\date{\today}

\begin{document}
\maketitle

\begin{abstract}
%% 6-10 sentences, written LAST. State: which runs (population, mass function, area), the
%% headline result per survey partition (Rubin / Roman / joint), and every caveat that changes
%% how a number should be read -- above all whether the runs predate a known bug (the
%% extinction law before e5ecb54 is the current example). Absolute yields only with their F.
\TODO{abstract}
\end{abstract}

% -------------------------------------------------------------------------------------------------
\section{Simulation and runs}
%% What the simulator does, in one paragraph (reuse populations_report.tex §1 almost verbatim),
%% then the runs. The table below is not optional: a reader must be able to tell a Monte Carlo
%% sample size from a yield, and the runs' commits from the code's current state.

\TODO{one paragraph on the simulator; the exact command line of each run}

\begin{table}[h]
\centering
\caption{Run configuration and outcome. ``Detected draws'' is a Monte Carlo \emph{sample size},
set by \texttt{-{}-events}/\texttt{-{}-lenses}, not a forecast; Section~\ref{sec:yield} converts it
to a yield.}
\label{tab:runs}
\begin{tabular}{lr}
\toprule
 & \TODO{run name} \\
\midrule
git commit of the run (\texttt{run\_provenance.txt}) & \TODO{} \\
population and mass function                          & \TODO{} \\
sightlines scanned / aggregated / barren              & \TODO{} \\
weightable draws                                      & \TODO{} \\
detected draws (MC)                                   & \TODO{} \\
$\Neff$ (weighted)                                    & \TODO{} \\
CPU time                                              & \TODO{} \\
\bottomrule
\end{tabular}
\end{table}

\paragraph{Weighting.}
%% Keep this paragraph in every report. The Monte Carlo lacks the rate's sqrt(M) v_t factor
%% (DEVIATIONS 41), so every pooled number is weighted and quoted with N_eff. An unweighted
%% number describes the sample, not the sky; say so explicitly if you ever show one.
Every pooled fraction, median or ratio carries the per-event weight
$W \propto w_{\rm area}\,N_\star/n_{\rm sim}\,\sqrt{\Ml}\,v_t\,Z(D_s)$ and is quoted with the
\citet{Kish1965} effective sample size $\Neff = (\sum w)^2/\sum w^2$.

% -------------------------------------------------------------------------------------------------
\section{\TODO{results section}}
%% One section per question. For each: figure, then a table of the numbers the figure shows,
%% then prose that says what the numbers MEAN (not what they are -- the table did that).
%% Rules the analysis layer enforces and the prose must respect:
%%   * -1.0 is the "not measured" sentinel; it never enters a sum, mean, ratio or histogram.
%%   * okA/okB (matrix inverted) is not "every parameter measured": check activePhotParams.
%%   * Rubin-only / Roman-only / joint curves each have their OWN denominator; say which.
%%   * sigma_joint <= sigma_single is a physical invariant; a violation is a bug, not a result.

\begin{figure}[h]
\centering
%% \includegraphics[width=\textwidth]{<figure>.png}
\TODO{figure}
\caption{\TODO{caption: what is plotted, over which sample, with what weighting}}
\label{fig:TODO}
\end{figure}

\TODO{what the figure shows and why}

% -------------------------------------------------------------------------------------------------
\section{Absolute event yields}
\label{sec:yield}
%% Include whenever a COUNT of events is quoted. Recompute with:
%%   .roman/bin/python analysis/y1_absolute_yield.py --run <name>=<run dir> ... -o <out>
%%   .roman/bin/python analysis/y3_yield_vs_F.py <out>/y1_yields.csv -o <out>/y3
%% and cite Report/populations/populations_report.tex §8 for the derivation rather than repeating it in full.
%% Always: (i) the F grid with its literature sources; (ii) which counting convention (per
%% resolved object, as the Sajadian papers use, or all stars); (iii) the tau check ratio,
%% which must be ~1 before any yield is trusted; (iv) N_1, the yield per unit F, so readers
%% can use their own F.

The expected number of events is $N(F) = F\,N_1$, with $N_1$ the yield at $F=1$
(derivation and validation in the compact-object report). Normalisation check, $\tau$ rebuilt
from the rate constants over the C++ optical depth: \TODO{pooled ratio per run}.

\begin{table}[h]
\centering
\caption{Yield per unit abundance, $N_1$, over the ten-year window, per-object convention.
\TODO{scope: whole scan or Roman footprint}}
\label{tab:n1}
\begin{tabular}{lr}
\toprule
selection & $N_1$ \\
\midrule
Roman detects                 & \TODO{} \\
Rubin detects                 & \TODO{} \\
both detect                   & \TODO{} \\
Roman, $\sigma(\Ml)/\Ml<10\%$ & \TODO{} \\
\bottomrule
\end{tabular}
\end{table}

% -------------------------------------------------------------------------------------------------
\section{Comparison with the literature}
%% Before comparing with any paper: align mass function, abundance, area, season count and
%% detection cut, and use THEIR denominator (detections vs characterised events). The first
%% SS23 comparison got two of these wrong and flattered the agreement (DEVIATIONS 54a).
\TODO{like-for-like comparison, or delete this section}

% -------------------------------------------------------------------------------------------------
\section{Caveats}
%% Copy the standing ones from OPEN_ITEMS.md that bear on these numbers; do not invent new
%% reassurance. At the time of writing the standing ones are: populations must not be pooled;
%% the model's latitude gradient is too shallow vs OGLE-IV; the footprint sample is small;
%% Rubin's astrometric error model is the weakest input; the parallax shortfall vs SS23.
\begin{itemize}
\item \TODO{does any number here predate a known bug fix? say which, and in which direction it biases}
\item \TODO{standing caveats from OPEN\_ITEMS.md}
\end{itemize}

% -------------------------------------------------------------------------------------------------
%% Reproduction block: the exact commands that made every figure and table. Keep it (commented
%% out if the reader does not need it) -- it is how the next session regenerates the report.
%\section*{Reproducing this report}
%\begin{verbatim}
%<analysis commands, one per figure/table>
%cd Report/<topic> && latexmk -pdf <topic>_report.tex
%\end{verbatim}

%% Shared bibliography. Add an entry to refs.bib only for a paper actually cited, and check its
%% journal reference against arXiv or the publisher -- never from memory.
\bibliographystyle{plainnat}
\bibliography{../refs}

\end{document}
